\documentclass[a4paper,12pt]{article}
\usepackage[T1]{fontenc}
\usepackage[utf8]{inputenc}
\usepackage[magyar]{babel}
\usepackage{palatino}
\usepackage{amsmath}   
\usepackage{amssymb}
% struktogram rajzolásához be kell tölteni a .tex forrásfájllal azonos könyvtárba helyezett stuki.sty-t
\usepackage{stuki}
% specifikációban címkézett változó írásához, A=(x:Z,l:L) alakban
\newcommand{\lv}{{:}}

% a specifikáció utófeltételében használt nagy függvények definiálása
\DeclareMathOperator*{\SEARCH}{\textit{\small \textbf{SEARCH}}}
\DeclareMathOperator*{\SELECT}{\textit{\small \textbf{SELECT}}}
\DeclareMathOperator*{\MAX}{\textit{\small \textbf{MAX}}}

\begin{document}
	\author{Zsolt Borsi}	
	\bigskip
	\bigskip

\begin{center}
{\Large Intervallumos programozási tételek}
\end{center}
\bigskip

%%% ÖSSZEGZÉS TÉTEL
\textbf{Összegzés}\\
\medskip
%%% specifikáció
$A = (m\lv\mathbb{Z}, n\lv\mathbb{Z}, s\lv\mathbb{Z})$\\[2pt]
$Ef = (m = m' \land n = n')$\\[2pt]
$Uf = (Ef \land s = \sum\limits_{i=m}^{n}f(i))$

%%% struktogram
\medskip
\begin{stuki}
  \stm{s:= 0}
  \begin{WHILE}{1}{\stm{i = m .. n}}
      \stm{s:= s + f(i)}
  \end{WHILE}
\end{stuki}
\bigskip

%%% SZÁMLÁLÁS TÉTEL
\textbf{Számlálás}\\
\medskip
%%% specifikáció
$A = (m\lv\mathbb{Z}, n\lv\mathbb{Z}, c\lv\mathbb{Z})$\\[2pt]
$E\!f = ( m=m' \land n=n' )$\\[6pt]
$U\!f = ( E\!f \land c = \sum\limits_{\substack{i=m\\[2pt]  \beta(i)}}^{n} 1)$
% A\substack{arg1\\ arg2} azért kell, mert mind az arg1 és az arg2 a nagyoperátor alá kell kerüljön, külön sorba

%%% struktogram
\medskip
\begin{stuki}
  \stm{c:= 0}
  \begin{WHILE}{2}{\stm{i = m..n}}
    \begin{IF}{1}{\stm{\beta(i)}}
      \stm{c:= c+1}
      \ELSE
      \stm{SKIP}
    \end{IF}
  \end{WHILE}
\end{stuki}

\bigskip

%%% MAXIMUMKIVÁLASZTÁS TÉTEL
\textbf{Maximumkiválasztás}\\
\medskip
%%% specifikáció
$A = (m\lv\mathbb{Z}, n\lv\mathbb{Z}, max\lv\mathcal{H}, ind\lv\mathbb{Z})$\\[2pt]
$E\!f = ( m=m' \land n=n' \land m \le n )$\\[6pt]
$U\!f = (E\!f \land (max,ind) = \MAX\limits_{i=m}^{n} f(i))$

%%% struktogram
\begin{stuki}
  \stm{max, ind:= f(m), m}
  \begin{WHILE}{2}{\stm{i = m+1..n}}
    \begin{IF}{1}{\stm{max < f(i)}}
      \stm{max,ind:= f(i),i}
      \ELSE
      \stm{SKIP}
    \end{IF}
  \end{WHILE}
\end{stuki}

\bigskip

%%% KIVÁLASZTÁS TÉTEL
\textbf{Kiválasztás}\\
\medskip
%%% specifikáció
$A = (m\lv\mathbb{Z}, i\lv\mathbb{Z})$\\[2pt]
$E\!f = ( m=m' \land \exists k \ge m\colon \beta(k) )$\\[6pt]
$U\!f = (E\!f \land i = \SELECT\limits_{i\ge m}^{} \beta(i))$

%%% struktogram
\medskip
\begin{stuki}[4cm]
  \stm{i:= m}
  \begin{WHILE}{1}{\stm{\neg\beta(i)}}
      \stm{i:= i + 1}
  \end{WHILE}
\end{stuki}
\bigskip

%%% LINEÁRIS KERESÉS TÉTEL
\textbf{Lineáris keresés}\\
\medskip
%%% specifikáció
$A = (m\lv\mathbb{Z}, n\lv\mathbb{Z}, l\lv\mathbb{L}, ind\lv\mathbb{Z})$\\[2pt]
$E\!f = ( m=m' \land n=n' )$\\[6pt]
$U\!f = (E\!f \land (l,ind) = \SEARCH\limits_{i=m}^{n} \beta(i))$

%%% struktogram
\medskip
\begin{stuki}
  \stm{l,i:= hamis,m}
  \begin{WHILE}{2}{\stm{\neg l \land i \le n}}
      \stm{l,ind:= \beta(i),i}
      \stm{i:= i + 1}
  \end{WHILE}
\end{stuki}
\bigskip

%%% OPTIMISTA ELDÖNTÉS TÉTEL
\textbf{Optimista eldöntés}\\
\medskip
%%% specifikáció
$A = (m\lv\mathbb{Z}, n\lv\mathbb{Z}, l\lv\mathbb{L})$\\[2pt]
$E\!f = (m=m' \land n=n')$\\[6pt]
$U\!f = (E\!f \land l = \forall\!\SEARCH\limits_{i=m}^{n} \beta(i))$

%%% struktogram
\medskip
\begin{stuki}
  \stm{l,i:= igaz,m}
  \begin{WHILE}{2}{\stm{l \land i \le n}}
      \stm{l:= \beta(i)}
      \stm{i:= i + 1}
  \end{WHILE}
\end{stuki}
\bigskip

%%% FELTÉTELES MAXIMUMKERESÉS TÉTEL
\textbf{Feltételes maximumkeresés}\\
\medskip
$A = (m\lv\mathbb{Z}, n\lv\mathbb{Z}, l\lv\mathbb{L}, max\lv\mathcal{H}, ind\lv\mathbb{Z})$\\[2pt]
$E\!f = (m=m' \land n=n')$\\[6pt]
$U\!f = (E\!f \land (max,ind) = \MAX\limits_{\substack{i=m\\[2pt] \beta(i)}}^{n} f(i))$

%%% struktogram
\medskip
\begin{stuki}[12cm]
  \stm{l:= hamis}
  \begin{WHILE}{4}{\stm{i = m..n}}
    \begin{CASE}{3}{16}
      \WHEN[3]{\stm{\neg\beta(i)}}
      \stm{SKIP}
      \WHEN[8]{\stm{l \land \beta(i)}}
      \begin{IF}{2}{\stm{max < f(i)}}
        \stm{max,ind:= \\f(i),i}
        \ELSE
        \stm{SKIP}
      \end{IF}
      \WHEN[5]{\stm{\neg l \land \beta(i)}}
        \stm{l,max,ind:= \\igaz,f(i),i}
    \end{CASE}
  \end{WHILE}
\end{stuki}

\end{document}